% !TeX TS-program = pdflatex
\documentclass[11pt,a4paper]{article}
\usepackage[margin=0.9in,headheight=14pt]{geometry}
\usepackage[T1]{fontenc}
\usepackage{lmodern}
\usepackage{amsmath,amssymb}
\usepackage{enumitem,booktabs,fancyhdr}
\setlength{\parindent}{0pt}
\setlength{\parskip}{0.5em}
\linespread{1.08}
\pagestyle{fancy}
\fancyhf{}
\fancyhead[C]{\small Probability \& Statistics for EECS: Homework \#04}
\fancyfoot[C]{\thepage}
\renewcommand{\headrulewidth}{0.4pt}
\renewcommand{\footrulewidth}{0.4pt}
\setlist[enumerate]{label=(\alph*),leftmargin=*,itemsep=1em}
\newcommand{\Prob}{P}
\title{SI 140A.02 Probability \& Statistics for EECS, Fall 2026\\[0.4em]Homework 4}
\author{Name: Jianxuan Yao \\ Student ID: 2025521024}
\date{
Due date: Due on Oct. 18, 2026, 23:59 UTC+8}
\begin{document}
\maketitle
\thispagestyle{plain}
\vspace{1em}
There are five problems in this assignment.

Read all the instructions below carefully before you start working on the assignment, and before you make a submission.
\begin{itemize}
\item You are required to write down all the major steps towards making your conclusions; otherwise you may receive limited credit for the problem.
\item Write your homework in English; otherwise you will receive no credit for this homework.
\item Any form of plagiarism will result in no credit for this homework.
\item You are required to submit your homework in PDF format generated by LaTeX, otherwise you will receive no credit for this homework.
\end{itemize}

\clearpage
\section*{Problem 1}
A token starts at position $1$ on the four positions $0,1,2,3$. It stops when it first reaches $0$ or $3$. Fix $r\in[0,1]$. All random choices in different rounds are independent. In an interior position, either of the following rules can be used throughout the process:
\begin{itemize}
\item \textbf{Pause rule:} With probability $r$, the token stays where it is; with probability $1-r$, a fresh fair coin moves it one position left or right.
\item \textbf{Restart rule:} With probability $r$, the token is placed at position $1$; with probability $1-r$, a fresh fair coin moves it one position left or right.
\end{itemize}
\begin{enumerate}
\item Under the pause rule, determine the probability of ever reaching position $3$ for every $r\in[0,1]$. Justify your answer.
\item Answer part (a) under the restart rule.
\item A student claims that pausing and restarting have the same effect because both only delay the token. Determine exactly when the two success probabilities are equal. For $r<1$, let $T$ be the number of rounds until the token stops. Find $\Prob(T>n)$ for every integer $n\geq0$ under each rule, and use your result to prove that each process stops with probability $1$. \emph{Hint:} Condition on either possible interior position after $n$ rounds.
\end{enumerate}

\textbf{Solution:}

\begin{enumerate}
    \item[(a)]

    When $r=1$, the token never moves, so the probability of reaching position $3$ is $0$.

    When $r<1$, let $p_k$ be the probability of eventually reaching position $3$ starting from position $k$. To position $k \in {1, 2}$, we have

    $$
    p_k = r \cdot p_k + \frac{1-r}{2} p_{k-1} + \frac{1-r}{2} p_{k+1}
    $$

    $$
    p_k = \frac{1}{2} p_{k-1} + \frac{1}{2} p_{k+1}
    $$

    And because of $p_0 = 0$ and $p_3 = 1$, 

    $$
    \begin{cases}
    p_1 = \frac{1}{2} p_0 + \frac{1}{2} p_2 = \frac{1}{2} p_2 \\
    p_2 = \frac{1}{2} p_1 + \frac{1}{2} p_3 = \frac{1}{2} p_1 + \frac{1}{2}
    \end{cases}
    $$

    Solving the system of equations, we get $p_1 = \frac{1}{3}$ and $p_2 = \frac{2}{3}$.
    
    Therefore, the probability of eventually reaching position $3$ starting from position $1$ is $\frac{1}{3}$ for all $r<1$.

    \item[(b)]

    When $r=1$, the token never moves, so the probability of reaching position $3$ is $0$.

    When $r<1$, let $p_k$ be the probability of eventually reaching position $3$ starting from position $k$. To position $k \in {1, 2}$, we have
    
    $$
    p_k = r p_1 + \frac{1-r}{2} p_{k-1} + \frac{1-r}{2} p_{k+1}
    $$

    $$
    \begin{cases}
    p_1 = r p_1 + \frac{1-r}{2} p_0 + \frac{1-r}{2} p_2 \\
    p_2 = r p_1 + \frac{1-r}{2} p_1 + \frac{1-r}{2} p_3
    \end{cases}
    $$

    And because of $p_0 = 0$ and $p_3 = 1$, we have

    $$
    \begin{cases}
        p_1 = \frac{1}{2} p_2 \\
        p_2 = \frac{1+r}{2} p_1 + \frac{1-r}{2}
    \end{cases}
    $$

    $$
    p_1 = \frac{1-r}{3-r}
    $$

    \item[(c)]
    
    $$
    \frac{1}{3} = \frac{1-r}{3-r}
    $$

    $$
    r = 0
    $$
    
    So $r = 0$ or $r = 1$.

    Let $a_n = P(X_n = 1, T > n)$ and $b_n = P(X_n = 2, T > n)$, we have $P(T > n) = a_n + b_n$.

    In Pause rule, we have

    $$
    \begin{pmatrix} a_{n+1} \\ b_{n+1} \end{pmatrix}
    =
    \begin{pmatrix}
    r & \frac{1-r}{2} \\
    \frac{1-r}{2} & r
    \end{pmatrix}
    \begin{pmatrix} a_n \\ b_n \end{pmatrix}
    $$

    $$
    P(T_{\text{pause}} > n) = a_n + b_n = \left(\frac{1+r}{2}\right)^n
    $$

    In Restart rule, we have

    $$
    \begin{pmatrix} a_{n+1} \\ b_{n+1} \end{pmatrix}
    =
    \begin{pmatrix}
    r & \frac{1+r}{2} \\
    \frac{1-r}{2} & 0 
    \end{pmatrix}
    \begin{pmatrix} a_n \\ b_n \end{pmatrix}
    $$

    $$
    \Prob(T_{\text{restart}} > n) = \left(\frac{1+r}{2}\right)^n
    $$

    So

    $$
    P(T = \infty)
    = \lim_{n \to \infty} P(T > n)
    = \lim_{n \to \infty} \left(\frac{1+r}{2}\right)^n
    = 0
    $$

    Therefore,

    $$
    P(T < \infty) = 1 - P(T = \infty) = 1
    $$
\end{enumerate}

\clearpage
\section*{Problem 2}
Two algorithms, $A$ and $B$, are each tested on $100$ tasks. Across the $200$ test records, exactly $100$ tasks are easy and $100$ are hard. Their exact success rates are:
\begin{center}
\begin{tabular}{ccc}
\toprule Algorithm & Easy tasks & Hard tasks\\\midrule
$A$ & $4/5$ & $2/5$\\
$B$ & $9/10$ & $3/5$\\\bottomrule
\end{tabular}
\end{center}
Each algorithm's easy/hard versus success/failure table has strictly positive integer entries. Let $x$ be the number of easy tasks assigned to $A$.
\begin{enumerate}
\item Find every feasible $x$ for which $A$ has a strictly greater overall success rate than $B$. For the smallest such $x$, give both complete $2\times2$ count tables.
\item From this part onward, use a probability model without finite-sample or integer-count constraints. The displayed within-type success probabilities remain fixed. Both algorithms receive the same easy-task proportion $t\in[0,1]$. An audit reports that their overall success probabilities differ by $13/100$, but omits which algorithm performs better. Recover $t$, identify the better algorithm, and calculate both overall success probabilities.
\item Now let $s,t\in(0,1)$ be the easy-task proportions for $A,B$, respectively, so each algorithm receives both task types. Construct one pair $(s,t)$ satisfying $|s-t|\leq2/5$ for which $A$'s overall success probability is exactly $1/100$ greater than $B$'s. Calculate both probabilities and check all the constraints for your pair. Briefly explain how this is compatible with $B$'s superiority within each task type.
\end{enumerate}

\textbf{Solution:}

\begin{enumerate}
    \item[(a)]

    $x$ must be a multiple of $10$ and $x \in [0,100]$.

    $$
    \frac{4}{5}x + \frac{2}{5}(100-x) > \frac{9}{10}(100-x) + \frac{3}{5}x
    $$

    $$
    x \geq \frac{500}{7}
    $$

    And considering the integer constraint, we have $x \in \{80, 90\}$.

    \begin{center}
    \begin{tabular}{ccc}
    \toprule A & Easy & Hard \\\midrule
    Success & 64 & 8 \\
    Failure & 16 & 12 \\\bottomrule
    \end{tabular}
    \end{center}

    \begin{center}
    \begin{tabular}{ccc}
    \toprule B & Easy & Hard \\\midrule
    Success & 18 & 48 \\
    Failure & 2 & 32 \\\bottomrule
    \end{tabular}
    \end{center}

    \item[(b)]
    
    $$
    P(A)=\frac{4}{5}t + \frac{2}{5}(1-t) = \frac{2}{5}t + \frac{2}{5}
    $$

    $$
    P(B)=\frac{9}{10}t + \frac{3}{5}(1-t) = \frac{3}{10}t + \frac{3}{5}
    $$

    $$
    \vert P(A) - P(B) \vert 
    = \vert \frac{1}{10}t - \frac{1}{5} \vert
    = \frac{13}{100}
    $$

    And $t \in [0,1]$, so $t = 0.7$.

    B is the better algorithm.

    $$
    P(A) = \frac{2}{5} \cdot 0.7 + \frac{2}{5} = \frac{17}{25}
    $$

    $$
    P(B) = \frac{3}{10} \cdot 0.7 + \frac{3}{5} = \frac{81}{100}
    $$

    \item[(c)]
    
    $$
    P(A) = \frac{2}{5}s + \frac{2}{5}
    $$

    $$
    P(B) = \frac{3}{10}t + \frac{3}{5}
    $$

    $$
    P(A) - P(B) = \frac{1}{100}
    $$

    $$
    \frac{2}{5}s - \frac{3}{10}t - \frac{1}{5} = \frac{1}{100}
    $$

    Select $s=0.9$ and $t=0.5$.

    Though $B$ is better in each task type, $A$ has a higher overall success probability because it receives a much higher proportion of easy tasks.
\end{enumerate}

\clearpage
\section*{Problem 3}
An automated station is placed in mode $H$ with probability $1/3$ and mode $L$ with probability $2/3$. The mode is chosen once for three checks. Let $A_i$ mean that check $i$ passes. Each check passes with probability $3/4$ conditional on $H$, and with probability $1/4$ conditional on $L$. Let $U=A_1\cap A_2\cap A_3$.
\begin{enumerate}
\item First assume the three pass events are mutually independent conditional on either mode. Compute $\Prob(H\mid U)$. Are $A_1$ and $A_2$ independent without conditioning on the mode? Support your answer with probabilities, and explain the difference between the two independence statements.
\item For the rest of this problem, assume only that every pair of pass events is independent conditional on each mode. For a fixed mode, write $p$ for the common pass probability and $z$ for the probability of three passes. Complete the table below and justify why the outcomes in each grouped row must have equal probabilities. Write $1$ for a pass and $0$ for a failure.
\begin{center}
\begin{tabular}{cc}
\toprule Outcome & Conditional probability of each outcome\\\midrule
$111$ & $z$\\
$110,101,011$ & $\cdots$\\
$100,010,001$ & $\cdots$\\
$000$ & $\cdots$\\\bottomrule
\end{tabular}
\end{center}
Find the full feasible intervals for $h=\Prob(U\mid H)$ and $\ell=\Prob(U\mid L)$. Explain why every pair $(h,\ell)$ in those intervals gives a valid model with the specified marginals and pairwise independence.
\item Find the exact set of possible values of $\Prob(H\mid U)$ under the assumptions of part (b). Prove both that no other value is possible and that every value in your set is attainable. Explain why the single answer in part (a) cannot be inferred from pairwise independence alone.
\end{enumerate}

\textbf{Solution:}

\begin{enumerate}
    \item[(a)]

    $$
    P(H \mid U)
    = \frac{P(U \mid H) P(H)}{P(U \mid H) P(H) + P(U \mid L) P(L)}
    = \frac{(\frac 34)^3 \cdot \frac 13}{(\frac 34)^3 \cdot \frac 13 + (\frac 14)^3 \cdot \frac 23}
    = \frac{27}{29}
    $$

    $$
    P(A_1)
    = P(A_1 \mid H) P(H) + P(A_1 \mid L) P(L)
    = \frac 34 \cdot \frac 13 + \frac 14 \cdot \frac 23
    = \frac{5}{12}
    $$

    Similarly, $P(A_2) = \frac{5}{12}$.

    $$
    P(A_1 \cap A_2)
    = P(A_1 \cap A_2 \mid H) P(H) + P(A_1 \cap A_2 \mid L) P(L)
    = (\frac 34)^2 \cdot \frac 13 + (\frac 14)^2 \cdot \frac 23
    = \frac{13}{48}
    $$

    $$
    P(A_1) P(A_2) = \frac{5}{12} \cdot \frac{5}{12} = \frac{25}{144} \neg \frac{13}{48}
    $$

    $A_1$ and $A_2$ are not independent without conditioning on the mode.

    The difference between the two independence statements is that conditional independence means that the events are independent given a certain condition, while unconditional independence means that the events are independent without any conditions. In this case, $A_1$ and $A_2$ are independent when we know the mode, but they are not independent when we do not know the mode.

    \item[(b)]
    
    $$
    P(110) = P(A_1 \cap A_2) - P(A_1 \cap A_2 \cap A_3) = p^2 - z
    $$

    Similarly, $P(101) = p^2 - z$ and $P(011) = p^2 - z$.

    $$
    P(100) = P(A_1) - P(A_1 \cap A_2) - P(A_1 \cap A_3) + P(A_1 \cap A_2 \cap A_3) = p - 2(p^2 - z) - z = p - 2p^2 + z
    $$

    $$
    P(010) = P(001) = p - 2p^2 + z
    $$

    $$
    P(000) = 1 - P(111) - 3P(110) - 3P(100) = 1 - z - 3(p^2 - z) - 3(p - 2p^2 + z) = 1 - 3p + 3p^2 - z
    $$

    \begin{center}
    \begin{tabular}{cc}
    \toprule Outcome & Conditional probability of each outcome\\\midrule
    $111$ & $z$\\
    $110,101,011$ & $p^2 - z$\\
    $100,010,001$ & $p - 2p^2 + z$\\
    $000$ & $1 - 3p + 3p^2 - z$\\\bottomrule
    \end{tabular}
    \end{center}

    $$
    \begin{cases}
        z \ge 0 \\
        p^2 - z \ge 0 \implies z \le p^2 \\
        p - 2p^2 + z \ge 0 \implies z \ge 2p^2 - p \\
        1 - 3p + 3p^2 - z \ge 0 \implies z \le 1 - 3p + 3p^2
    \end{cases}
    $$

    $$
    \max(0, 2p^2 - p) \le z \le \min(p^2, 1 - 3p + 3p^2)
    $$

    For mode H, $h \in \left[\frac{3}{8}, \frac{7}{16}\right]$.
    
    For mode L, $l \in \left[0, \frac{1}{16}\right]$.

    The conditional probabilities of mode H and mode L are independently constructed in two disjoint probability spaces. $h$ and $l$ can be selected independently in their respective intervals.

    \item[(c)]

    $$
    P(H \mid U)
    = \frac{P(U \mid H) P(H)}{P(U \mid H) P(H) + P(U \mid L) P(L)}
    = \frac{h \cdot \frac{1}{3}}{h \cdot \frac{1}{3} + l \cdot \frac{2}{3}}
    = \frac{h}{h + 2l}
    = \frac{1}{1 + 2\frac{l}{h}}
    $$

    $$
    \frac{l}{h} \in \left[0, \frac{1}{6}\right]
    $$

    So $\Prob(H \mid U) \in \left[\frac{3}{4}, 1\right]$.

    According to the monotonicity of $f(x) = \frac{1}{1 + 2x}$, any value outside of $\left[\frac{3}{4}, 1\right]$ is impossible.

    For any $\alpha \in \left[\frac{3}{4}, 1\right]$, the ratio $u^\star = \frac{1-\alpha}{2\alpha} \in [0, \frac{1}{6}]$. Fix $h = \frac{3}{8}$, then $l = \frac{3}{8} u^\star$.

    $0 \le u^* \le \frac{1}{6}$, so $l \in [0, \frac{1}{16}]$. Then we prove every value of $P(H \mid U)$ can be achieved.

    Pairwise independence only imposes constraints on the probabilities of second-order intersections $P(A_i \cap A_j) = p^2$, but it completely leaves unconstrained the probability of third-order intersections $P(A_1 \cap A_2 \cap A_3) = z$.
\end{enumerate}

\clearpage
\section*{Problem 4}
In a three-door game, the car is behind doors $1,2,3$ with probabilities
$1/2,1/3,1/6$. You choose door $1$. The host knows the car's location and
opens a goat door from doors $2,3$. If both doors contain goats, host A
opens door $2$ with probability $1/4$, whereas host B does so with
probability $3/4$. Otherwise, the opening is forced.

A fair coin selects one host, independently of the car. The selected host
first performs a calibration round, with the car known to be behind door
$1$ and door $1$ initially selected, and then the main game. Given the host, the two rounds use independent
random choices. Let $C$ and $J$ be the calibration and main openings,
respectively, and let $K$ be the main car door.

System F retains exactly the runs with $C=2$; system M retains exactly
the runs with $C=J$. Write $F=\{C=2\}$ and $M=\{C=J\}$. In a retained
run, the system and $J$ are reported, but the host and $C$ are not
separately reported. All questions below concern the main game.

\begin{enumerate}
\item Prove that the two systems give the same
posterior probability of $K=1$ when $J=2$ is reported, and calculate
this probability. Identify the conditioning events before calculating.

\item For each system, make a two-row table, indexed
by $j=2,3$, giving
\[
\Prob(J=j\mid\text{retained}),\qquad
\Prob(K=1\mid J=j,\text{retained}).
\]
Reuse your result from part (a). Apply the conditional law of total
probability to find \mbox{$\Prob(K=1\mid F)$} and \mbox{$\Prob(K=1\mid M)$}.
Explain why the equality in part (a) does not imply equality of these
two probabilities.

\item Modify only the host selection: now select a
host independently by a fair coin for each of the two rounds. Conditional on the selected hosts, the rounds use independent random choices. Both
recording systems keep their original rules. Does either system change
the joint distribution of $(K,J)$ by retaining runs? Prove your answer
for every pair $(k,j)$ with positive probability. Begin by finding the
probability of retention conditional on $K=k$ and $J=j$, using the law of
total probability over the calibration host. Explain what feature of
the original experiment allowed the two systems to give different
answers in part (b).
\end{enumerate}

\textbf{Solution:}

\begin{enumerate}
    \item[(a)]

    When $J=2$, $M = \{C=J\} = \{C=2\} = F$, so the two system give the same information, and they share the same posterior probability.

    $$
    \begin{aligned}
    &P(K=1 \mid J=2, F) \\
    =& P(K=1 \mid C=2, J=2) \\
    =& \frac{P(C=2, J=2 \mid K=1) P(K=1)}{P(C=2, J=2 \mid K=1) P(K=1) + P(C=2, J=2 \mid K=3) P(K=3)} \\
    =& \frac{P(C=2, J=2 \mid K=1) \cdot \frac{1}{2}}{P(C=2, J=2 \mid K=1) \cdot \frac{1}{2} + P(C=2, J=2 \mid K=3) \cdot \frac{1}{6}}
    \end{aligned}
    $$

    $$
    P(C=2, J=2 \mid K=1) = P(C=2, J=2, A \mid K=1) + P(C=2, J=2, B \mid K=1)
    $$

    $$
    \begin{aligned}
    P(C=2, J=2, A \mid K=1)
    &= P(C=2, J=2 \mid A, K=1) P(A \mid K=1) \\
    &= P(C=2, J=2 \mid A, K=1) P(A) \\
    &= \left(\frac{1}{4}\right)^2 \cdot \frac{1}{2} \\
    &= \frac{1}{32}
    \end{aligned}
    $$

    $$
    P(C=2, J=2, B \mid K=1)
    = P(C=2, J=2 \mid B, K=1) P(B)
    = \left(\frac{3}{4}\right)^2 \cdot \frac{1}{2}
    = \frac{9}{32}
    $$

    $$
    P(C=2, J=2 \mid K=1) = \frac{5}{16}
    $$

    $$
    P(C=2, J=2 \mid K=3) = P(C=2, J=2, A \mid K=3) + P(C=2, J=2, B \mid K=3)
    $$

    $$
    \begin{aligned}
    P(C=2, J=2, A \mid K=3)
    &= P(C=2, J=2 \mid A, K=3) P(A \mid K=3) \\
    &= P(C=2, J=2 \mid A, K=3) P(A) \\
    &= P(C=2 \mid A, K=3) P(J=2 \mid C=2, A, K=3) P(A) \\
    &= \frac{1}{4} \cdot \frac{1}{2} \\
    &= \frac{1}{8}
    \end{aligned}
    $$

    $$
    P(C=2, J=2, B \mid K=3) = \frac{3}{8}
    $$

    $$
    P(C=2, J=2 \mid K=3) = \frac{1}{8} + \frac{3}{8} = \frac{1}{2}
    $$

    $$
    P(K=1 \mid J=2, F) = \frac{\frac{5}{16} \cdot \frac{1}{2}}{\frac{5}{16} \cdot \frac{1}{2} + \frac{1}{2} \cdot \frac{1}{6}} = \frac{15}{23}
    $$

    \item[(b)]
    
    $$
    P(C=2) = P(H=A)P(C=2 \mid H=A) + P(H=B)P(C=2 \mid H=B) = \frac{1}{2} \cdot \frac{1}{4} + \frac{1}{2} \cdot \frac{3}{4} = \frac{1}{2}
    $$

    $$
    P(J=2 \mid F) = \frac{P(C=2, J=2)}{P(C=2)} = \frac{\frac{23}{96}}{\frac{1}{2}} = \frac{23}{48}
    $$

    $$
    P(J=3 \mid F) = 1 - P(J=2 \mid F) = 1 - \frac{23}{48} = \frac{25}{48}
    $$

    $$
    \begin{aligned}
    &P(C=2, J=3 \mid K=1) \\
    =& \frac{1}{2}\left(P(C=2 \mid A)P(J=3 \mid K=1, A) + P(C=2 \mid B)P(J=3 \mid K=1, B)\right) \\
    =& \frac{1}{2}\left(\frac{1}{4} \times \frac{3}{4} + \frac{3}{4} \times \frac{1}{4}\right) \\
    =& \frac{1}{2} \left(\frac{3}{16} + \frac{3}{16}\right) \\
    =& \frac{3}{16}
    \end{aligned}
    $$

    $$
    P(K=1, C=2, J=3)
    = P(K=1)P(C=2, J=3 \mid K=1)
    = \frac{1}{2} \times \frac{3}{16} = \frac{3}{32}
    = \frac{9}{96}
    $$

    $$
    P(K=1 \mid J=3, F)
    = \frac{P(K=1, C=2, J=3)}{P(C=2, J=3)}
    = \frac{\frac{9}{96}}{\frac{25}{96}}
    = \frac{9}{25}
    $$

    \begin{center}
    \begin{tabular}{ccc}
    \toprule $j$ & $P(J=j \mid F)$ & $P(K=1 \mid J=j, F)$ \\\midrule
    $2$ & $\frac{23}{48}$ & $\frac{15}{23}$ \\
    $3$ & $\frac{25}{48}$ & $\frac{9}{25}$ \\\bottomrule
    \end{tabular}
    \end{center}

    $$
    \begin{aligned}
    &P(C=3, J=3 \mid K=1) \\
    =& \frac{1}{2}\left(P(C=3 \mid A)P(J=3 \mid K=1, A) + P(C=3 \mid B)P(J=3 \mid K=1, B)\right) \\
    =& \frac{1}{2}\left(\frac{3}{4} \times \frac{3}{4} + \frac{1}{4} \times \frac{1}{4}\right) \\
    =& \frac{5}{16}
    \end{aligned}
    $$

    $$
    P(C=3, J=3 \mid K=2) = P(C=3) = \frac{1}{2}
    $$

    $$
    P(J=3, M) = P(C=3, J=3) = \frac{15}{96} + \frac{16}{96} = \frac{31}{96}
    $$

    $$
    P(M) = P(J=2, M) + P(J=3, M) = \frac{23}{96} + \frac{31}{96} = \frac{54}{96} = \frac{9}{16}
    $$

    $$
    P(J=2 \mid M) = \frac{P(J=2, M)}{P(M)} = \frac{\frac{23}{96}}{\frac{54}{96}} = \frac{23}{54}
    $$

    $$
    P(J=3 \mid M) = \frac{P(J=3, M)}{P(M)} = \frac{\frac{31}{96}}{\frac{54}{96}} = \frac{31}{54}
    $$

    $$
    P(K=1 \mid J=3, M) = \frac{P(K=1, C=3, J=3)}{P(C=3, J=3)} = \frac{\frac{15}{96}}{\frac{31}{96}} = \frac{15}{31}
    $$

    \begin{center}
    \begin{tabular}{ccc}
    \toprule $j$ & $P(J=j \mid M)$ & $P(K=1 \mid J=j, M)$ \\\midrule
    $2$ & $\frac{23}{54}$ & $\frac{15}{23}$ \\
    $3$ & $\frac{31}{54}$ & $\frac{15}{31}$ \\\bottomrule
    \end{tabular}
    \end{center} 

    $$
    \begin{aligned}
    P(K=1 \mid F)
    &= P(J=2 \mid F) P(K=1 \mid J=2, F) + P(J=3 \mid F) P(K=1 \mid J=3, F) \\
    &= \left(\frac{23}{48} \times \frac{15}{23}\right) + \left(\frac{25}{48} \times \frac{9}{25}\right) \\
    &= \frac{15}{48} + \frac{9}{48} \\
    &= \frac{24}{48} \\
    &= \frac{1}{2}
    \end{aligned}
    $$

    $$
    \begin{aligned}
    P(K=1 \mid M)
    &= P(J=2 \mid M) P(K=1 \mid J=2, M) + P(J=3 \mid M) P(K=1 \mid J=3, M) \\
    &= \left(\frac{23}{54} \times \frac{15}{23}\right) + \left(\frac{31}{54} \times \frac{15}{31}\right) \\
    &= \frac{15}{54} + \frac{15}{54} \\
    &= \frac{30}{54} \\
    &= \frac{5}{9}
    \end{aligned}
    $$

    $$
    P(K=1 \mid J=3, F) = \frac{9}{25} \neq \frac{15}{31} = P(K=1 \mid J=3, M)
    $$

    And the probability distributions of $J$ are different in the two systems, so the equality in part (a) does not imply equality of these two probabilities.

    \item[(c)]
    
    $$
    \begin{aligned}
    P(F \mid K=k, J=j)
    &= P(C=2 \mid K=k, J=j) \\
    &= \sum_{h \in \{A, B\}} P(H_C = h \mid K=k, J=j) P(C=2 \mid H_C = h, K=k, J=j) \\
    &= \sum_{h \in \{A, B\}} P(H_C = h) P(C=2 \mid H_C = h) \\
    &= \frac{1}{2} \cdot P(C=2 \mid H_C=A) + \frac{1}{2} \cdot P(C=2 \mid H_C=B) \\
    &= \frac{1}{2} \left(\frac{1}{4}\right) + \frac{1}{2} \left(\frac{3}{4}\right) \\
    &= \frac{1}{8} + \frac{3}{8} \\
    &= \frac{1}{2}
    \end{aligned}
    $$

    $$
    P(K=k, J=j \mid F) 
    = \frac{P(F \mid K=k, J=j) P(K=k, J=j)}{P(F)}
    = \frac{\frac{1}{2} \cdot P(K=k, J=j)}{\frac{1}{2}}
    = P(K=k, J=j)
    $$

    F does not change the joint distribution of $(K,J)$ by retaining runs.

    $$
    \begin{aligned}
    P(M \mid K=k, J=j) &= P(C=j \mid K=k, J=j) \\
    &= \sum_{h \in \{A, B\}} P(H_C = h \mid K=k, J=j) P(C=j \mid H_C = h, K=k, J=j) \\
    &= \sum_{h \in \{A, B\}} P(H_C = h) P(C=j \mid H_C = h)
    \end{aligned}
    $$

    $$
    P(C=2)
    = \frac{1}{2}P(C=2 \mid H_C=A) + \frac{1}{2}P(C=2 \mid H_C=B)
    = \frac{1}{2} \left(\frac{1}{4}\right) + \frac{1}{2} \left(\frac{3}{4}\right)
    = \frac{1}{2}
    $$

    $$
    P(C=3)
    = \frac{1}{2}P(C=3 \mid H_C=A) + \frac{1}{2}P(C=3 \mid H_C=B)
    = \frac{1}{2} \left(\frac{3}{4}\right) + \frac{1}{2} \left(\frac{1}{4}\right)
    = \frac{1}{2}
    $$

    $$
    P(K=k, J=j \mid M) = \frac{P(M \mid K=k, J=j) P(K=k, J=j)}{P(M)} = \frac{\frac{1}{2} \cdot P(K=k, J=j)}{\frac{1}{2}} = P(K=k, J=j)
    $$

    M does not change the joint distribution of $(K,J)$ by retaining runs.

    In (b), the two systems give different answers because the two rounds share the same host, so the calibration round provides information about the main round.
    
    In (c), the two rounds have independent hosts, so the calibration round provides no information about the main round.
\end{enumerate}

\clearpage
\section*{Problem 5}
Start with one amoeba at time $0$. At each minute's end, every living
amoeba dies or splits into two. The death probability is $1/4$ in odd
minutes and $p\in(0,1)$ in even minutes. All choices are independent
across individuals and minutes; an empty population stays empty.
Do not use general branching-process theorems.

\begin{enumerate}
\item For integers $n\geq0$, let $x_n,y_n$ be the probabilities of
extinction within $n$ updates, starting with one amoeba under the
alternating rules, with the odd/even rule first, respectively. Let $x,y$ be the corresponding
eventual probabilities. Find the initial values, first-step recursions
and limiting equations. Explain why $x=y$ cannot be assumed.
You may use \mbox{$\Prob(\bigcup_{n\geq0}E_n)=\lim_{n\to\infty}\Prob(E_n)$}
whenever $E_n\subseteq E_{n+1}$ for all $n\geq0$.

\item Set $p=2/3$. Prove $x=y=1$. An averaged model replaces both
death probabilities by their average $\bar p=11/24$, keeping all other rules unchanged;
its eventual extinction probability is given as $11/13$.
Calculate extinction probabilities by minute $2$ in both models.
Does averaging overestimate or underestimate extinction at minute $2$
and eventually?

\item Now $p$ is unknown, and the probability of exactly two amoebae
at the end of minute $2$ is $1/3$. Given exactly four at the end of
minute $3$, find every possible probability that all four descend
from the same one of the two minute-$1$ amoebae, with the corresponding
$p$. Is this conditional probability uniquely determined?
\end{enumerate}

\textbf{Solution:}

\begin{enumerate}
    \item[(a)]

    $$
    x_0 = 0, \quad y_0 = 0
    $$

    $$
    x_{n+1} = \frac{1}{4} + \frac{3}{4} y_n^2
    $$

    $$
    y_{n+1} = p + (1-p) x_n^2
    $$

    $$
    x = \lim_{n \to \infty} x_n, \quad y = \lim_{n \to \infty} y_n
    $$

    $$
    \begin{cases}
    x = \dfrac{1}{4} + \dfrac{3}{4} y^2 \\
    y = p + (1-p) x^2
    \end{cases}
    $$

    The state transitions are not symmetric, so $x=y$ cannot be assumed.

    \item[(b)]
    
    $$
    \begin{cases}    
    x = \dfrac{1}{4} + \dfrac{3}{4} y^2 \\
    y = \dfrac{2}{3} + \dfrac{1}{3} x^2
    \end{cases}
    $$

    $$
    x
    = \frac{1}{4} + \frac{3}{4} \left(\frac{2}{3} + \frac{1}{3} x^2\right)^2
    = \frac{1}{4} + \frac{3}{4} \left(\frac{4}{9} + \frac{4}{9} x^2 + \frac{1}{9} x^4\right)
    = \frac{1}{4} + \frac{1}{3} + \frac{1}{3} x^2 + \frac{1}{12} x^4
    $$

    $$
    (x - 1)^2 (x^2 + 2x + 7) = 0
    $$

    $$
    x = 1
    $$

    $$
    y = \frac{2}{3} + \frac{1}{3} x^2 = 1
    $$

    Origin:

    $$
    x_1 = \frac{1}{4} + \frac{3}{4} y_0^2 = \frac{1}{4}
    $$

    $$
    y_1 = \frac{2}{3} + \frac{1}{3} x_0^2 = \frac{2}{3}
    $$

    $$
    x_2 = \frac{1}{4} + \frac{3}{4} y_1^2 = \frac{7}{12}
    $$

    Average:

    $$
    z_1 = \bar{p} = \frac{11}{24}
    $$

    $$
    z_2 = \bar{p} + (1-\bar{p}) z_1^2 = \frac{7909}{13824}
    $$

    $$
    x_2 = \frac{7}{12} = \frac{8064}{13824} > \frac{7909}{13824} = z_2
    $$

    The averaged model underestimates extinction at minute $2$.

    The eventual extinction probability of the averaged model is $11/13 < 1$, so it also underestimates eventual extinction.

    \item[(c)]
    
    $$
    P(Z_2 = 2 \mid Z_1 = 2) = 2 \cdot p(1-p)
    $$

    $$
    P(Z_2 = 2) = P(Z_1 = 2) \cdot P(Z_2 = 2 \mid Z_1 = 2) =  \frac{3}{2} p(1-p)
    $$

    $$
    p = \frac{1}{3} \text{ or } p = \frac{2}{3}
    $$

    Let $A$ be $\{Z_3 = 4\}$ and $B$ be the four amoebae descend from the same one of the two minute amoebae.

    Let the two amoebae surviving to minute $2$ be ancestor $I$ and $II$. Let $W_1$ and $W_2$ be the descendant number of them respectively.

    For one amoeba, after two minutes, the probability of $0$ is $p$ and the probability of $2$ is $1-p$.

    And after three minutes, for one amoeba surviving to minute $2$,

    $$
    P(W = 0) = p + (1-p) \cdot \left(\frac{1}{4}\right)^2 = p + \frac{1-p}{16}
    $$

    $$
    P(W = 2) = (1-p) \cdot \left[2 \cdot \frac{1}{4} \cdot \frac{3}{4}\right] = \frac{6(1-p)}{16} = \frac{3(1-p)}{8}
    $$

    $$
    P(W = 4) = (1-p) \cdot \left(\frac{3}{4}\right)^2 = \frac{9(1-p)}{16}
    $$

    $$
    P(A, B \mid Z_1 = 2)
    = 2 \cdot P(W = 4) \cdot P(W = 0)
    = 2 \cdot \left[\frac{9(1-p)}{16}\right] \cdot \left[p + \frac{1-p}{16}\right]
    = \frac{9(1-p)}{8} \left(\frac{15p + 1}{16}\right)
    $$

    $$
    P(W_1 = 2, W_2 = 2 \mid Z_1 = 2)
    = [P(W = 2)]^2 = \left[\frac{6(1-p)}{16}\right]^2
    = \frac{9(1-p)^2}{64}
    $$

    $$
    \begin{aligned}
    P(A \mid Z_1 = 2)
    &= P(A, B \mid Z_1 = 2) + P(W_1 = 2, W_2 = 2 \mid Z_1 = 2) \\
    &= \frac{9(1-p)}{128} \left[2(15p + 1) + 2(1-p)\right] \\
    &= \frac{9(1-p)}{128} (30p + 2 + 2 - 2p) \\
    &= \frac{9(1-p)(28p + 4)}{128} \\
    &= \frac{9(1-p)(7p + 1)}{32}
    \end{aligned}
    $$

    $$
    P(B \mid A)
    = \frac{P(B \cap A \mid Z_1 = 2)}{P(A \mid Z_1 = 2)}
    = \frac{15p + 1}{4(7p + 1)}
    $$

    When $p=\frac{1}{3}$, $P(B \mid A) = \frac{9}{20}$.

    When $p=\frac{2}{3}$, $P(B \mid A) = \frac{33}{68}$.

    It is not uniquely determined.
\end{enumerate}

\end{document}
